\chapter{FEASIBILITY ANALYSIS}

\section{Technical Feasibility}
 The proposed automated hydroponic nutrient dosing and monitoring system is technically feasible due to the use of readily available, well-documented hardware components and established control methodologies. The ESP32 microcontroller provides sufficient processing capability, communication interfaces, and integrated wireless connectivity to support real-time sensing, closed-loop control, and data logging. The adopted threshold-based proportional control strategy with stabilization delay is computationally lightweight and well suited to the slow dynamic behavior of hydroponic nutrient solutions, avoiding the complexity and tuning challenges of advanced controllers while maintaining stability and reliability. Similar system architectures have been successfully implemented in prior IoT-based hydroponic studies, confirming the practicality of the proposed design.
 
\section{Economic Feasibility}
 From an economic perspective, the project is feasible and cost-effective for small-scale and academic applications. All major components, including sensors, dosing pumps, and the control unit, are low-cost and widely available, resulting in a total system cost significantly lower than commercial hydroponic controllers. In addition to low initial investment, the system is expected to reduce operational costs by minimizing nutrient wastage through precise dosing and reducing labor requirements via automation. The use of open-source hardware and software further eliminates licensing costs, enhancing affordability and scalability.
 
\section{Operational Feasibility}
 Operationally, the system is designed for autonomous and reliable operation with minimal human intervention. Continuous monitoring and automated corrective actions ensure stable nutrient conditions without requiring frequent manual adjustments. Built-in safety features, stabilization delays, and calibration routines enhance robustness and reduce the risk of overdosing or sensor-related errors. The availability of a user interface for real-time monitoring and historical data review enables easy supervision, making the system practical for users with limited technical expertise and suitable for continuous operation in small-scale hydroponic environments.